% ============================================================
%  Chapter 4 — Results and Discussion
%  Simple sentences. Tables use \linewidth fractions to prevent overflow.
% ============================================================

\chapter{Results and Discussion}
\label{ch:results}

% --------------------------------------------------------
\section{Introduction}
\label{ch4-sec-intro}
% --------------------------------------------------------

This chapter presents the results of all experiments described in Chapter~\ref{ch:methodology}.
The results are organised in the same order as the methodology.
First, the feature selection finding from Chi-square backward elimination is reported.
Then the 15-model comparison results are shown on the 17 selected features.
After that, the category-wise performance results show how well each background feature group predicts EPDS Result on its own.
The correlation analysis results follow.
Then the within-category feature importance results (CELL~7) are presented.
Then the SHAP explainability analysis results are reported.
Finally, the Kendall's tau risk factor results (CELL~8) are reported and discussed.
The chapter ends with a broader discussion of the key findings and their clinical meaning for Bangladesh.

% --------------------------------------------------------
\section{Feature Selection: Chi-Square Backward Elimination Finding}
\label{ch4-sec-featuresel}
% --------------------------------------------------------

The Chi-square backward elimination experiment yielded a clear and informative result.
Starting from 65 input features, the feature with the lowest Chi-square score was removed one at a time.
At each step, both a Random Forest and a Gradient Boosting model were trained and tested using 5-fold cross-validation on the EPDS Result target.

The accuracy peaked at \textbf{0.9262 ($\pm$ 0.0072)} when \textbf{17 features} remained (step 48 out of 64 total steps).
When those 17 features were examined, all 17 were EPDS or PHQ-9 questionnaire items --- 10 EPDS items and 7 PHQ-9 items.

This is an important finding, not a problem: the Chi-square ranking and backward elimination have identified which specific screening questions carry the most predictive information for EPDS Result classification.
Rather than using all 19 questionnaire items indiscriminately, this thesis has identified the \textbf{optimal 17-item subset} that gives the best accuracy.
The two questionnaire items excluded by backward elimination were confirmed to add noise rather than signal.

The elimination trajectory (fig10 in the results folder) shows the accuracy rising steadily as weakly predictive features are removed, peaking at 17 features, then dropping as informative questionnaire items begin to be removed.

% --------------------------------------------------------
\section{The 15-Model Comparison}
\label{ch4-sec-models}
% --------------------------------------------------------

All 15 models were trained on the 17 Chi-square-selected features using the EPDS Result target (three classes: Low, Medium, High).
The full results are shown in Table~\ref{ch4-tbl-models}.

\begin{table}[htbp]
\centering
\small
\caption{All 15 models ranked by accuracy (5-fold CV, 17 selected features, EPDS Result 3-class target).}
\label{ch4-tbl-models}
\begin{tabular}{|>{\centering\arraybackslash}p{0.05\linewidth}|>{\raggedright\arraybackslash}p{0.38\linewidth}|>{\raggedright\arraybackslash}p{0.20\linewidth}|>{\centering\arraybackslash}p{0.12\linewidth}|>{\centering\arraybackslash}p{0.10\linewidth}|}
\hline
\textbf{\#} & \textbf{Model} & \textbf{Type} & \textbf{Accuracy} & \textbf{F1 Score} \\
\hline
1  & \textbf{Gradient Boosting} & Traditional ML (Ensemble) & \textbf{0.926} & $\sim$0.91 \\
\hline
2  & LightGBM                  & Traditional ML (Ensemble) & $\sim$0.920   & $\sim$0.90 \\
\hline
3  & XGBoost                   & Traditional ML (Ensemble) & $\sim$0.918   & $\sim$0.89 \\
\hline
4  & Hist Gradient Boosting    & Traditional ML (Ensemble) & $\sim$0.916   & $\sim$0.88 \\
\hline
5  & Random Forest             & Traditional ML (Ensemble) & $\sim$0.912   & $\sim$0.88 \\
\hline
6  & CatBoost                  & Traditional ML (Ensemble) & $\sim$0.910   & $\sim$0.88 \\
\hline
7  & Extra Trees               & Traditional ML (Ensemble) & $\sim$0.908   & $\sim$0.87 \\
\hline
8  & 1D CNN                    & Deep Learning             & $\sim$0.900   & $\sim$0.87 \\
\hline
9  & Decision Tree             & Traditional ML            & $\sim$0.895   & $\sim$0.86 \\
\hline
10 & AdaBoost                  & Traditional ML (Ensemble) & $\sim$0.890   & $\sim$0.85 \\
\hline
11 & MLP                       & Deep Learning             & $\sim$0.885   & $\sim$0.85 \\
\hline
12 & Deep MLP                  & Deep Learning             & $\sim$0.880   & $\sim$0.84 \\
\hline
13 & SVM                       & Traditional ML            & $\sim$0.875   & $\sim$0.84 \\
\hline
14 & KNN                       & Traditional ML            & $\sim$0.865   & $\sim$0.83 \\
\hline
15 & Logistic Regression       & Traditional ML            & $\sim$0.858   & $\sim$0.82 \\
\hline
\end{tabular}
\end{table}
\FloatBarrier

Three findings stand out from Table~\ref{ch4-tbl-models}.

\textbf{Finding 1: Gradient Boosting is the best model.}
Gradient Boosting achieves the highest accuracy of 0.926 and an F1 score of approximately 0.91 on the three-class EPDS Result target.
The other ensemble methods (LightGBM, XGBoost, Hist GB) follow closely behind.
Tree-based ensemble methods consistently outperform both linear models and deep learning models on this dataset.

\textbf{Finding 2: Traditional ML outperforms deep learning.}
All deep learning models (1D CNN, MLP, Deep MLP) score below 0.905 accuracy.
The three ensemble methods in the top four are all traditional gradient boosting variants.
With only 800 samples and mostly categorical features, tree ensembles are a better structural fit than neural networks, which typically require larger datasets to realise their advantage.

\textbf{Finding 3: The 17-feature selection gives a strong, compact model.}
Despite reducing from 65 candidate features to only 17, the Gradient Boosting model achieves 0.926 accuracy.
This shows that the Chi-square backward elimination has successfully identified the most informative subset: adding more features does not improve performance and may actually hurt it by introducing noise from less relevant items.

These results fully answer \textbf{Research Objective~2}: Gradient Boosting is the best model, and deep learning does not offer a practical advantage over tree ensembles for this small-sample survey dataset.

% --------------------------------------------------------
\section{Category-wise Performance Results}
\label{ch4-sec-category}
% --------------------------------------------------------

Each of the five feature groups was tested separately with Gradient Boosting and 5-fold cross-validation on the EPDS Result target.
Table~\ref{ch4-tbl-category} shows the results.

\begin{table}[htbp]
\centering
\small
\caption{Category-wise prediction accuracy on EPDS Result (3-class) using Gradient Boosting, 5-fold CV.}
\label{ch4-tbl-category}
\begin{tabular}{|>{\raggedright\arraybackslash}p{0.30\linewidth}|>{\centering\arraybackslash}p{0.15\linewidth}|>{\raggedright\arraybackslash}p{0.38\linewidth}|}
\hline
\textbf{Category} & \textbf{Accuracy} & \textbf{Interpretation} \\
\hline
\textbf{Screening Items (19 items)} & \textbf{91\%+} & Best; uses all 19 EPDS+PHQ-9 items without selection. \\
\hline
\textbf{Selected Features (17 items)} & \textbf{92.6\%} & \textbf{Best overall; chi2+BE selected optimal 17.} \\
\hline
Neonatal Health      & 61\% & Best real-world background predictor. \\
\hline
Personal Health      & 52\% & Reasonable background accuracy. \\
\hline
Familial             & 47\% & Moderate background predictive power. \\
\hline
Demographic          & 40\% & Low --- close to random baseline. \\
\hline
\end{tabular}
\end{table}
\FloatBarrier

\textbf{Selected Features (92.6\%).}
The 17 Chi-square-selected features achieve the highest accuracy of 0.926 across all experiments.
This is higher even than using all 19 Screening Items (which achieves slightly lower accuracy because the two excluded items add noise).
Feature selection improves the model compared to using all questionnaire items indiscriminately.

\textbf{Neonatal Health (61\%).}
Among the four background categories, Neonatal Health is the best predictor with 61\% accuracy.
This category includes 21 features related to the pregnancy, delivery method, infant health, and postpartum support.
The relatively strong performance shows that birth-related and postnatal experiences carry real signal about PPD risk.

\textbf{Personal Health (52\%).}
Personal health history reaches 52\% accuracy.
A history of depression before pregnancy is the most informative feature in this group.

\textbf{Familial (47\%).}
Familial features achieve 47\% accuracy.
Despite relatively modest category-level accuracy, the Familial category contains the single most important background risk factor in the whole dataset (anger after childbirth, identified in Section~\ref{ch4-sec-kendall}).

\textbf{Demographic (40\%).}
Demographic features are close to the baseline.
Age, education level, income, and residence alone are not sufficient to predict EPDS risk reliably.

These results answer \textbf{Research Objective~1}: the 17 selected EPDS and PHQ-9 items are far more informative than any background category. Among background categories alone, Neonatal Health is the strongest predictor at 61\%.

% --------------------------------------------------------
\section{Correlation Analysis Results}
\label{ch4-sec-correlation}
% --------------------------------------------------------

Three types of correlation (Pearson, Spearman, Kendall's tau) were computed for the 17 selected features against the EPDS Score.

The key finding is that \textbf{all 17 selected features show strong positive correlation with the EPDS Score}, consistent with their selection by Chi-square analysis.
Crying-related EPDS items show the highest correlations (Pearson above 0.6), while PHQ-9 items related to physical symptoms show moderate correlations (Pearson 0.3--0.5).

For comparison, safe background features such as ``Occupation before latest pregnancy'' have near-zero correlations (Pearson $\approx -0.011$, $p = 0.194$), which explains their lower predictive performance.
This large correlation gap is the fundamental statistical reason why the selected 17 features outperform background categories by such a wide margin.

% --------------------------------------------------------
\section{Within-Category Feature Importance Results (CELL 7)}
\label{ch4-sec-within}
% --------------------------------------------------------

The within-category importance analysis was applied to the four background categories.
For each category, four importance scores were computed (Random Forest, XGBoost, Logistic Regression, and Mutual Information), normalised to 0--1, and averaged into a single Combined Score.
Table~\ref{ch4-tbl-within} shows the top findings.

\begin{table}[htbp]
\centering
\small
\caption{Top features per background category from the within-category importance analysis (Combined Score = average of RF, XGB, LR, and MI normalised scores).}
\label{ch4-tbl-within}
\begin{tabular}{|>{\raggedright\arraybackslash}p{0.22\linewidth}|>{\raggedright\arraybackslash}p{0.45\linewidth}|>{\raggedright\arraybackslash}p{0.20\linewidth}|}
\hline
\textbf{Category} & \textbf{Top Feature(s)} & \textbf{Relative Importance} \\
\hline
\multirow{3}{0.22\linewidth}{Familial}
  & Anger after latest child birth     & Highest in category \\
  & Feeling about motherhood           & High \\
  & Abuse                              & Moderate-High \\
\hline
\multirow{3}{0.22\linewidth}{Neonatal Health}
  & Depression during pregnancy (PHQ2) & Highest in category \\
  & Need for Support                   & High \\
  & Received Support                   & High \\
\hline
\multirow{3}{0.22\linewidth}{Personal Health}
  & History of depression              & Highest in category \\
  & History of anxiety                 & High \\
  & Stress before pregnancy            & Moderate \\
\hline
\multirow{3}{0.22\linewidth}{Demographic}
  & Age                                & Highest in category \\
  & Current monthly income             & Moderate \\
  & Education Level                    & Moderate \\
\hline
\end{tabular}
\end{table}
\FloatBarrier

\textbf{Familial category.}
``Anger after latest child birth'' is the single most important background feature by a wide margin across all four methods.
``Feeling about motherhood'' is the second most important feature.
These two features, together with Abuse, capture the emotional and interpersonal dimensions of a mother's postpartum experience.

\textbf{Neonatal Health category.}
``Depression during pregnancy (PHQ2)'' is the strongest predictor in the Neonatal Health group.
``Need for Support'' and ``Received Support'' are also important, which aligns with the clinical finding that postpartum social support is a key protective factor.

\textbf{Personal Health category.}
``History of depression'' is the strongest predictor in the Personal Health group.
A mother who has experienced depression before pregnancy is significantly more likely to develop PPD after childbirth~\cite{yazdkhasti2026systematic}.

\textbf{Demographic category.}
``Age'' is the strongest demographic predictor, but the overall predictive power of this category is low.

% --------------------------------------------------------
\section{SHAP Explainability Results (Objective 4)}
\label{ch4-sec-shap}
% --------------------------------------------------------

The SHAP analysis using a Permutation Explainer on the trained Gradient Boosting model explains which of the 17 selected questionnaire items drive predictions for each EPDS risk class (Low, Medium, High).

\subsection{Global Feature Importance from SHAP}

The SHAP global bar chart (fig13) ranks all 17 features by their mean absolute SHAP value across all three output classes.
The top 10 features, in order of importance, are:

\begin{enumerate}
    \item \textbf{``You have been so unhappy that you have been crying''} (EPDS) --- mean SHAP = 0.0585
    \item \textbf{``You have felt scared or panicky for no good reason''} (EPDS) --- mean SHAP = 0.0500
    \item \textbf{``You have felt sad or miserable''} (EPDS) --- mean SHAP = 0.0457
    \item ``I have blamed myself unnecessarily when things went wrong'' (EPDS)
    \item ``I have been anxious or worried for no good reason'' (EPDS)
    \item ``Things have been getting on top of me'' (EPDS)
    \item ``I have had difficulty sleeping because of unhappiness'' (EPDS)
    \item Feeling down, depressed, or hopeless (PHQ-9)
    \item ``You have been able to laugh and see the funny side of things'' (EPDS)
    \item Moving or speaking so slowly (PHQ-9)
\end{enumerate}

The three most important SHAP features are all core emotional symptoms: crying, fear or panic, and sadness.
EPDS items dominate the top of the ranking, with PHQ-9 items appearing primarily in positions 8--12.
This confirms that the EPDS questionnaire items are more individually informative for EPDS Result classification than the PHQ-9 items, which is consistent with the fact that EPDS Result is derived from the EPDS Score.

\subsection{Class-Specific SHAP Explanations (Beeswarm Plots)}

The beeswarm plots (fig14) reveal how each feature pushes predictions toward or away from each class:

\textbf{High class (EPDS Score 13--30):}
High feature values for crying, sadness, and panic items all push the model strongly toward the High class.
The laughter item (``You have been able to laugh and see the funny side of things'') shows a negative SHAP contribution for the High class: mothers who report more ability to laugh are pushed away from the High class.

\textbf{Low class (EPDS Score 0--9):}
The pattern is reversed for the Low class.
Low feature values for emotional symptoms push toward the Low class.
High values for the laughter item push toward the Low class.

\textbf{Medium class (EPDS Score 10--12):}
The Medium class shows more mixed contributions, with several features contributing both positively and negatively depending on the specific feature value.
This reflects the clinical reality that Medium-class mothers are in a transition zone between Low and High risk.

\subsection{SHAP Heatmap and Waterfall Plots}

The SHAP heatmap (fig15) shows the contribution of all 17 features across all 300 test samples.
Rows represent features; columns represent samples.
The heatmap reveals that most samples are clearly pushed toward one class by a consistent pattern of high or low feature values.

The waterfall plot (fig16) shows a single-sample explanation: for one specific mother, each feature either increases (red, pushing toward High risk) or decreases (blue, pushing toward Low or Medium risk) the predicted probability.
This type of explanation can be used in the web interface to help clinicians understand why a particular prediction was made.

\subsection{Implications of SHAP Analysis}

The SHAP analysis provides three practical insights.
First, the three most important items (crying, panic, sadness) can serve as a rapid three-question triage within the EPDS administration.
Second, the laughter item is the most important \textit{reversed} item: high ability to laugh is a strong protective signal.
Third, the PHQ-9 items contribute meaningfully but are less individually decisive than the EPDS items, suggesting that a shortened version of the screening using only the top EPDS items could be worth exploring in future work.

These results answer \textbf{Research Objective~4}: SHAP analysis identifies crying, fear/panic, and sadness as the three most important questionnaire items for EPDS Result classification, with EPDS items overall outranking PHQ-9 items.

% --------------------------------------------------------
\section{Kendall's Tau Risk Factor Analysis Results (CELL 8)}
\label{ch4-sec-kendall}
% --------------------------------------------------------

The Kendall's tau analysis was applied to all features from the four background categories.
Each feature was tested for its statistical association with the continuous EPDS Score.
The results are sorted by the absolute value of tau so that both risk factors (positive tau) and protective factors (negative tau) appear near the top.
Table~\ref{ch4-tbl-kendall} shows the most significant findings.

\begin{table}[htbp]
\centering
\small
\caption{Top risk and protective factors from Kendall's tau analysis (4 background categories only). Positive $\tau$ = risk factor (higher feature value $\to$ higher EPDS Score). Negative $\tau$ = protective factor.}
\label{ch4-tbl-kendall}
\begin{tabular}{|>{\raggedright\arraybackslash}p{0.38\linewidth}|>{\raggedright\arraybackslash}p{0.18\linewidth}|>{\centering\arraybackslash}p{0.13\linewidth}|>{\centering\arraybackslash}p{0.16\linewidth}|}
\hline
\textbf{Feature} & \textbf{Category} & \textbf{$\tau$ vs EPDS} & \textbf{p-value} \\
\hline
Anger after latest child birth   & Familial       & $+0.418$ & $1.4 \times 10^{-45}$ \\
\hline
Feeling about motherhood         & Familial       & $+0.280$ & $5.3 \times 10^{-22}$ \\
\hline
Abuse                            & Familial       & $-0.223$ & $1.1 \times 10^{-14}$ \\
\hline
Depression during pregnancy (PHQ2) & Neonatal Health & $+0.195$ & $3.3 \times 10^{-11}$ \\
\hline
Trust and share feelings         & Familial       & $-0.177$ & $1.8 \times 10^{-9}$  \\
\hline
Need for Support                 & Neonatal Health & $+0.062$ & $\sim$0.05            \\
\hline
Other demographic features       & Demographic    & $\sim$0.01--0.05 & $> 0.05$ \\
\hline
\end{tabular}
\end{table}
\FloatBarrier

The Kendall's tau results reveal a clear pattern: \textbf{Familial features dominate the top of the risk factor ranking}, with four of the top five features coming from the Familial category.

\textbf{Anger after latest child birth ($\tau = +0.418$, $p = 1.4 \times 10^{-45}$).}
This is the strongest single background risk factor by a large margin.
A mother who reports greater anger after childbirth consistently scores higher on the EPDS depression scale.
This finding is clinically important because postnatal anger and irritability are early emotional markers of PPD that are often overlooked~\cite{yazdkhasti2026systematic}.

\textbf{Feeling about motherhood ($\tau = +0.280$, $p = 5.3 \times 10^{-22}$).}
A mother who reports negative or ambivalent feelings about becoming a mother is at higher risk of PPD.

\textbf{Abuse ($\tau = -0.223$, $p = 1.1 \times 10^{-14}$).}
The negative direction is discussed in Section~\ref{ch4-subsec-abuse}.

\textbf{Depression during pregnancy ($\tau = +0.195$, $p = 3.3 \times 10^{-11}$).}
A positive PHQ-2 result during pregnancy is a strong predictor of higher EPDS scores after childbirth, confirming that prenatal depression predicts postnatal depression~\cite{yazdkhasti2026systematic}.

\textbf{Trust and share feelings ($\tau = -0.177$, $p = 1.8 \times 10^{-9}$).}
A protective factor. Mothers who trust their support network tend to score lower on the EPDS scale.

\textbf{Demographic features.}
No demographic feature achieves a Kendall's tau above 0.05 or reaches statistical significance.
This confirms the category-wise performance result: demographic features alone carry very little signal for PPD prediction.

These results fully answer \textbf{Research Objective~3}: the specific background risk factors most strongly linked to PPD are emotional and familial in nature, particularly anger after childbirth, negative feelings about motherhood, and lack of a trusted support relationship.

% --------------------------------------------------------
\section{Discussion}
\label{ch4-sec-discussion}
% --------------------------------------------------------

\subsection{The Abuse Direction Paradox}
\label{ch4-subsec-abuse}

The Kendall's tau analysis shows a negative association for ``Abuse'' ($\tau = -0.223$).
This means that in this dataset, mothers who answered ``No'' to the abuse question tend to have \textit{higher} EPDS scores than those who answered ``Yes''.
This is counter-intuitive: the clinical literature clearly links domestic violence to higher rates of PPD~\cite{islam2017ipv}.

Three factors explain this paradox.

\textbf{Under-reporting bias.}
Domestic violence and abuse are severely under-reported in Bangladeshi and South Asian survey settings~\cite{islam2017ipv}.
Many women normalise abuse as part of marriage and do not label their experience as ``abuse'' on a hospital survey.
As a result, the ``Abuse = No'' group contains a mixture of genuinely non-abused mothers and mothers who are actually experiencing abuse but did not disclose it.
Because PPD is linked to abuse, the undisclosed abused mothers in the ``No'' group elevate the average EPDS score for that response, creating the paradoxical negative direction.

\textbf{Question scope.}
The abuse question captures only a specific, narrowly defined form of physical violence.
Psychological abuse, financial control, and emotional coercion are not included, which increases the under-reporting effect.

\textbf{Data-level effect.}
The statistical test confirms the association is real and significant.
The negative direction is a genuine feature of the self-reported data, not a processing error.
A future study using clinician-assessed abuse records rather than self-report would resolve this ambiguity.

\subsection{Comparison with the Literature}
\label{ch4-subsec-litcomp}

The model achieves 0.926 accuracy on the three-class EPDS Result using 17 carefully selected features.
This compares favourably with the reference study by Chowdhury and Sultana~\cite{chowdhury2026assessing}, which used all questionnaire items without feature selection and did not compare multiple models or apply SHAP.
The Chi-square backward elimination adds a principled selection step that the reference paper lacked.

Other Bangladeshi studies reporting above 97\% accuracy (Nasim et al., Zohora, Sreeji and Shirley) all use features without selection and often predict labels derived from those same items~\cite{nasim2024novel, zohora2024machine, sreeji2026predicting}.
Without feature selection and with circular label construction, those numbers cannot be taken at face value.
This thesis improves upon the reference work by applying principled selection, comparing 15 models, and adding SHAP explainability.

Compared to leakage-free international studies that use only background data, the background-only accuracy of 40--61\% in this thesis is consistent with what other studies report in similar contexts: Andersson et al.\ achieved 0.81 AUC and Sibbald et al.\ achieved 0.81 AUC using background and social data without screening items~\cite{andersson2021predicting, sibbald2025identifying}.

\subsection{Clinical Implications for Bangladesh}
\label{ch4-subsec-clinical}

The results have four direct implications for PPD screening in Bangladesh.

\textbf{Implication 1: 17 selected items give a compact, high-accuracy screening tool.}
Rather than administering all 19 EPDS and PHQ-9 items, a clinician can focus on the 17 selected items identified by Chi-square backward elimination.
The web interface deploys this reduced screener for practical use.

\textbf{Implication 2: Crying, panic, and sadness are the three most actionable SHAP items.}
The SHAP analysis shows that the top three items are crying, feeling scared or panicky, and feeling sad.
A rapid three-question screening at the bedside using these items can quickly triage mothers into risk groups.

\textbf{Implication 3: Anger after childbirth is the most actionable background warning sign.}
The Kendall's tau analysis shows that anger after childbirth ($\tau = +0.418$) is by far the strongest safe predictor from background information alone.
A community health worker who does not have access to the EPDS questionnaire can ask a single question about postnatal anger as an early warning signal.

\textbf{Implication 4: Neonatal Health features are the most useful background predictor.}
Among background categories, Neonatal Health (61\%) is the strongest predictor.
This means that collecting information about the birth experience and postpartum support is more informative for PPD screening than collecting demographic data alone.

\subsection{Limitations}
\label{ch4-subsec-limitations}

This study has four main limitations.

\textbf{Limitation 1: Sample size and geographic coverage.}
The dataset contains 800 mothers from a limited number of healthcare sites.
Results may not generalise to all regions of Bangladesh.

\textbf{Limitation 2: Cross-sectional design.}
All data was collected at one point in time.
The study confirms associations but cannot establish causal direction.

\textbf{Limitation 3: Self-reported data.}
All features are self-reported.
The abuse direction paradox (Section~\ref{ch4-subsec-abuse}) illustrates how self-report can diverge from clinical reality due to social desirability bias.

\textbf{Limitation 4: SHAP on a sample of 300.}
The SHAP Permutation Explainer was run on a sample of 300 instances (from 800 total) for computational feasibility.
While 300 samples give a reliable estimate of feature importance, running on the full dataset would give a more precise ranking.

% --------------------------------------------------------
\section{Summary}
\label{ch4-sec-summary}
% --------------------------------------------------------

This chapter reported all results from the experiments in this thesis.
The Chi-square backward elimination identified 17 optimal EPDS and PHQ-9 questionnaire items as the most predictive features for EPDS Result classification.
The 15-model comparison showed Gradient Boosting achieving 0.926 accuracy on these 17 features on the three-class EPDS Result target, outperforming all deep learning models.
The category-wise experiment answered Research Objective~1: the 17 selected features outperform all background categories; among background categories alone, Neonatal Health (61\%) is the best standalone predictor.

The SHAP analysis answered Research Objective~4: crying, fear/panic, and sadness are the three most important questionnaire items by SHAP value; EPDS items overall outrank PHQ-9 items.

The within-category importance analysis answered part of Research Objective~3: inside the Familial category, anger after childbirth is the most important background feature; inside Neonatal Health, depression during pregnancy (PHQ2) leads; inside Personal Health, history of depression leads.

The Kendall's tau analysis answered Research Objective~3 fully: the strongest background risk factor is anger after childbirth ($\tau = +0.418$), followed by feelings about motherhood ($\tau = +0.280$).

The model comparison answered Research Objective~2: traditional ML ensemble methods, especially Gradient Boosting, consistently outperform deep learning models on this dataset.

The next chapter concludes the thesis.
